\documentclass[oneside,12pt]{article}          % ----- article

\input{../../../Latex/inputs/loading_packages.tex}			% ----- Packages utiles
\input{../../../Latex/inputs/preambule_exos.tex}         	% ----- Préambule feuilles de TD
\input{../../../Latex/inputs/mymacros.tex}                 % ----- Macros mathématiques
\input{../../../Latex/inputs/mise_en_page.tex} 		   	% ----- Macros de mise en page
\input{../../../Latex/inputs/macros_tikz.tex} 		   	% ----- Macros utiles TikZ

\fancyhead[R]{T\up{ale} Technologique}            % En-tête droit
\fancyhead[C]{Suites arithmétiques}


%%%%%%%%%%%%%%%%%%%%%%%%%%%%%%%%%%%%%%%%%%%%%%%%%%%%%%%%%%%%%%%%%%%%%%%%%%%%%%%%%%%%%%%%%%%%%%%%%%%%%%%%%%%%%
%%%%%%%%%%%%%%%%%%%%%%%%%%%%%%%%%%%%%%%%%%%%%%%%%%%%%%%%%%%%%%%%%%%%%%%%%%%%%%%%%%%%%%%%%%%%%%%%%%%%%%%%%%%%%
%%%%%%%%%%%%%%%%%%%%%%%%%%%%%%%%%%%%%%%%%%%%%%%%%%%%%%%%%%%%%%%%%%%%%%%%%%%%%%%%%%%%%%%%%%%%%%%%%%%%%%%%%%%%%


\begin{document}

\section{Rappels sur les suites}

\subsection{Calculs explicites}
Calculer les $5$ premiers termes des suites numériques définies ci-après :

\begin{tabular}{p{0.3\textwidth} p{0.3\textwidth} p{0.3\textwidth}}
	$u_n = n^2 - 3n$ & $v_n = \dfrac{1-n}{1+2n}$ & $w_n = 2^n + n$
\end{tabular}



\subsection{Calculs récurrents}
Calculer les $5$ premiers termes des suites numériques définies ci-après :

\begin{tabular}{p{0.3\textwidth} p{0.3\textwidth} p{0.3\textwidth}}
	$\begin{cases} u_0 = 1 \\ u_{n+1} = 2u_n - 3\end{cases}$ &
	$\begin{cases} v_1 = 2 \\ v_{n+1} = 1 + \dfrac{2}{v_n} \end{cases}$ &
	$\begin{cases} w_0 = 1 \\ w_{n+1} = w_n + 2n \end{cases}$
\end{tabular}



\subsection{Modélisation}
Pour chacune des situations suivantes, proposer une suite définie par récurrence qui permet de la modéliser :
\begin{enumerate}
    \item Dans une tube a essai on place 5 bactéries. Le nombre de bactéries double à chaque heure.
    \item Un gisement de pétrole a produit 60000 barils en l'an 2000 et sa production augmente de 1500 barils chaque année.
    \item Pour la naissance de leur fille, M et Mme Rakotobe placent 1000\euro\, sur un compte bancaire à taux d'intérêts annuels de $5\%$
    \item Dans un jeu vidéo une guilde se forme avec 25 membres. Chaque mois la guilde parvient à recruter 10 nouveaux membres, mais 20\% des anciens membres arrêtent de jouer car ils préfèrent utiliser leur temps libre pour résoudre des problèmes de mathématiques.
\end{enumerate}



\subsection{Extrapolation}
Donner une formule de récurrence ou une expression explicite pour les suites suivantes :

\renewcommand{\arraystretch}{1.6}
\begin{tabular}{p{0.18\textwidth} p{0.15\textwidth} p{0.15\textwidth} p{0.15\textwidth} p{0.15\textwidth}}
    a) $u_0 = 2$ & $u_1 = 4$ & $u_2 = 6$ & $u_3 = 8$ & $u_4 = 10$ \\
    b) $u_0 = 4 $ & $ u_1 = 8$ & $u_2 = 16$ & $u_3 = 32$ & $u_4 = 64$ \\
    c) $u_1 = 1$ & $u_2 = 2$ & $u_3 = 4$ & $u_4 = 7$ & $u_5 = 11$ \\
    d) $u_2 = \frac{1}{2}$ & $u_3 = \frac{2}{3}$ & $u_4 = \frac{3}{4}$ & $u_5 = \frac{4}{5}$ & $u_6 = \frac{5}{6}$ 
\end{tabular}



\section{Suites arithmétiques}

\subsection{Premiers termes}
\begin{enumerate}
	\item On considère $(u_n)$ une suite arithmétique de premier terme $u_0 = 3$ et de raison $r = 2$ \\
	Calculer $u_1$ , $u_2$ et $u_3$
	\item On considère $(v_n)$ une suite arithmétique de premier terme $v_0 = -5$ et de raison $r = \dfrac{3}{2}$ \\
	Calculer $v_1$ , $v_2$ et $v_3$
\end{enumerate}



\subsection{Question de nature}
Les suites suivantes sont-elles des suites arithmétiques ? Vous justifierez votre réponse :

\hspace{2cm} $u_n = 3n + 5 \quad$ ; \quad $v_n = n^2 + n \quad$ ; \quad $w_n = \dfrac{n}{n+1} \quad$ ; \quad $x_n = (n+1)^2 - n^2$



\subsection{Progression arithmétique}
\begin{enumerate}
	\item Les trois nombres suivants sont-ils en progression arithmétique ? \quad $\dfrac{1}{3} \quad$ ; \quad $\dfrac{1}{4}$ \quad et \quad $\dfrac{1}{5}$
	\item Les trois nombres suivants sont-ils en progression arithmétique ? \quad $-5 \quad$ ; \quad $2$ \quad et \quad $9$
	\item Les trois nombres suivants sont-ils en progression arithmétique ? \quad $3 \quad$ ; \quad $\dfrac{23}{6}$ \quad et \quad $\dfrac{14}{3}$
\end{enumerate}



\subsection{Forme explicite}
\begin{enumerate}
	\item On considère $(u_n)$ une suite arithmétique de raison $r = -2$ et de premier terme $u_0 = 35$
	\begin{enumerate}
		\item Exprimer $u_n$ en fonction de $n$
		\item Calculer alors $u_{10}$ et $u_{100}$
	\end{enumerate}
	\item On considère $(v_n)$ une suite arithmétique de raison $r = \dfrac{2}{3}$ et de premier terme $v_0 = 2$
	\begin{enumerate}
		\item Exprimer $v_n$ en fonction de $n$
		\item Calculer alors $v_{54}$ et $v_{123}$
	\end{enumerate}
\end{enumerate}



\subsection{Petits pots de miel}
Une entreprise d'apiculture note ses productions de miel année après année. Les résultats collectés sont donnés dans le tableau ci-contre :
\begin{center}
	\begin{tabular}{| P{2cm} | P{2cm} | P{2cm} | P{2cm} | P{2cm} | P{2cm} | P{2cm} |} \hline
		Année & 2020 & 2021 & 2022 & 2023 & 2024 & 2025 \\ \hline
		Nombre de pots produits & $155\,100$  & $162\,450$ & $169\,050$ & $175\,900$ & $183\,150$ & $190\,500$ \\ \hline
	\end{tabular}
\end{center}
\begin{enumerate}
	\item Une suite arithmétique semble-t-elle adaptée pour modéliser l'évolution de la production de l'entreprise d'apiculture depuis l'année 2020 ?
	\item On note $u_n$ la nombre de pots vendus pas l'entreprise lors de l'année 2020 + $n$ (en milliers de pots) \\
	Quelle est la valeur de $u_0$ ?
	\item On admet que la suite $(u_n)$ est arithmétique de raison $r = 7$
	\begin{enumerate}
		\item Exprimer $u_n$ en fonction de $n$
		\item Avec ce modèle de prévision, combien de pots l'entreprise vendra-t-elle en 2035 ?
	\end{enumerate}
\end{enumerate}



\subsection{Elles ont perdu la raison}
\begin{enumerate}
	\item On considère une suite arithmétique $(u_n)$ vérifiant $u_8 = 19$ et $u_{14} = 57$
	\begin{enumerate}
		\item Déterminer $r$ la raison de la suite $(u_n)$
		\item En déduire alors la valeur de $u_0$
		\item Calculer la valeur du terme $u_{100}$
	\end{enumerate}
	\item On considère une suite arithmétique $(v_n)$ vérifiant $u_{23} = 48$ et $u_{31} = 34$
	\begin{enumerate}
		\item Déterminer $r$ la raison de la suite $(v_n)$
		\item En déduire alors la valeur de $v_0$
		\item Calculer la valeur du terme $u_{50}$
	\end{enumerate}
\end{enumerate}



\subsection{Problème de seuil}
\begin{enumerate}
	\item On considère une suite arithmétique $(u_n)$ de premier terme $u_0 = 5$ et de raison $r = 2$ \\
	À partir de quel rang $n$ la suite $(u_n)$ dépassera-t-elle la valeur $1\,000$ ?
	\item On considère une suite arithmétique $(v_n)$ de premier terme $v_0 = 521$ et de raison $r = -4$ \\
	À partir de quel rang $n$ la suite $(v_n)$ deviendra-t-elle négative ?
\end{enumerate}



\subsection{Repeupler la montagne}
En juin 2020, une population de 200 marmottes a été introduite dans un massif montagneux où cette espèce était absente. Un zoologue en charge de ce projet souhaite modéliser l’évolution de cette population en fonction du temps.

Il constate qu’entre juin 2020 et juin 2021, la population a augmenté de 20 individus. Le zoologue propose alors un modèle où la population augmente de 20 individus tous les ans.

On note alors $u_n$ la population de marmottes que l’on peut estimer à l’aide de ce modèle lors du mois de juin de l'année 2020 + $n$.
\begin{enumerate}
	\item Quelle est la nature de la suite $(u_n)$ ? Préciser sa raison $r$ et son premier terme $u_0$
	\item Exprimer $u_n$ en fonction de $n$
	\item À combien d'individus peut-on estimer la population de marmottes au mois de juin de l'année 2030 ?
	\item Le zoologue a décidé que lorsque la population de marmottes dépasserait 700 individus, il dépêcherait une équipe sur les lieux pour étudier leur assimilation dans le massif montagneux. D'après le modèle choisi, en quelle année devra-t-il prévoir son étude ?
\end{enumerate}



\subsection{Achat revente}
En 2025, Issa a acheté un téléphone portable coûtant $1\,100$\euro.\\
On admettra que la valeur de ce téléphone diminue de $120$\euro\, chaque année.

On note $u_n$ la valeur du téléphone de Issa en l'année 2025 + $n$.
\begin{enumerate}
	\item Déterminer la valeur de $u_0$ , $u_1$ et $u_2$ puis interpréter ces résultats dans le contexte de l'exercice.
	\item Exprimer $u_{n+1}$ en fonction de $u_n$ et en déduire la nature de la suite $(u_n)$
	\item Exprimer $u_n$ en fonction de $n$
	\item À partir de quelle année le téléphone de Issa aura-t-il perdu la moitié de sa valeur initiale ?
\end{enumerate}





\section{Sommes de termes}




\end{document}